% Generated draft: requires mathematical and source review.
\documentclass[11pt]{article}
\usepackage[margin=25mm]{geometry}
\usepackage{fontspec}
\setmainfont{Latin Modern Roman}
\usepackage{amsmath,amssymb,mathtools,unicode-math}
\setmathfont{Latin Modern Math}
\usepackage{enumitem,needspace,xurl,hyperref}
\setlength{\parindent}{0pt}
\setlength{\parskip}{5pt}
\setlength{\emergencystretch}{3em}
\newcommand{\N}{\mathbb{N}}
\newcommand{\Z}{\mathbb{Z}}
\newcommand{\Q}{\mathbb{Q}}
\newcommand{\R}{\mathbb{R}}
\newcommand{\C}{\mathbb{C}}
\newcommand{\F}{\mathbb{F}}
\newcommand{\A}{\mathbb{A}}
\newcommand{\PP}{\mathbb{P}}
\newcommand{\E}{\mathbb{E}}
\newcommand{\eps}{\varepsilon}
\newcommand{\dd}{\,\mathrm{d}}
\begin{document}
\section*{UnsolvedMath ID: 30002815}
\textbf{Source problem code:} catalogue entry
\subsection{Definitions and mathematical statement}
For a planar polygonal tessellation \(T\) of a sheet domain \(\Omega\subseteq\R^2\), a rigid folded realization, in the usual idealized interpretation, is a continuous map \(f:\Omega\to\R^3\) whose restriction to each face is a Euclidean isometry onto a planar polygon. Thus faces remain rigid, and adjacent faces remain attached along their common edges. A rigid folding is a continuous family of such realizations. One degree of freedom ordinarily means that the configuration space, after identifying configurations differing by an ambient rigid motion, is locally one-dimensional along the relevant folding branch. These conventions explain the terminology; the supplied record does not specify the admissible sheets, tessellations, folding branches, or collision conditions.

\textbf{Source question.} Given an arbitrary surface with intrinsic curvature, does there exist a Miura-like tessellation of the plane whose folded realization approximates that surface? Can such a tessellation additionally be made rigidly foldable with one degree of freedom?

A possible quantitative structure for the first question is
\[
\text{for every }S\in\mathcal S\text{ and every }\varepsilon>0,\quad
\text{there exist }T\in\mathcal M\text{ and a folded realization }f
\text{ with }\operatorname{err}(f,S)<\varepsilon.
\]
Here \(\mathcal S\), \(\mathcal M\), and \(\operatorname{err}\) are placeholders for, respectively, the intended class of curved surfaces, the intended class of Miura-like tessellations, and the intended approximation criterion. They are not defined by the source record. The additional mechanism question asks whether the approximating realization can belong to a rigid folding mechanism with one degree of freedom. It is not specified whether the mechanism must start at a planar configuration or whether one-degree-of-freedom behavior must hold globally, locally near the approximating configuration, or away from singular configurations.

\textbf{Unresolved qualifications.} The record does not specify surface regularity, topology, compactness, boundary conditions, an ambient realization, or whether approximation is intrinsic or extrinsic. It does not define ``Miura-like,'' specify whether the sheet is finite or the entire plane, or state whether cuts, holes, self-intersections, or self-contact are permitted. Consequently, the displayed quantifier pattern records a possible precision framework and is not a fully determined conjecture. The supplied dated assessment mentions rigidity obstructions for positive Gaussian curvature but gives no hypotheses or precise obstruction statement; no such theorem is asserted here.
\subsection{Short English statement}
The question asks whether every surface with intrinsic curvature can be approximated by a folded Miura-like planar tessellation, and whether such approximations can also have a rigid folding mechanism with one degree of freedom. The surface class, tessellation class, approximation criterion, and admissible folding motions remain unspecified. The supplied partially solved label and curvature-obstruction discussion are dated curation requiring independent review.
\subsection{Sources}
\begin{itemize}
\item Local UnsolvedMath registry, supplied record with permanent ID \(30002815\), title ``Rigid Miura Tessellations Approximating Curved Surfaces,'' exact source problem code \texttt{catalogue entry}. The record belongs to the Oberwolfach Reports Open Problems collection and gives proposed year 2015; no individual proposer is supplied.
\item Supplied report attribution: \emph{Discrete Differential Geometry} (2015), Oberwolfach Reports, DOI 10.4171/owr/2015/13, pp.\ 661--729. \url{https://doi.org/10.4171/owr/2015/13}. This is also the supplied literature reference labelled ``OWR Report 13/2015.'' The record describes Oberwolfach Reports as documenting mathematical workshops and their open-problem discussions; it does not identify a problem-specific page or speaker.
\item Supplied literature reference: \emph{Curved display based on origami tessellations}. \url{https://doi.org/10.1038/s41378-021-00319-1}.
\item Literature assessment dated exactly \texttt{2026-08-21}: the supplied status is \texttt{partially\_solved}. The supplied assessment says that origami constructions approximate broad classes of curved surfaces, that universal Miura-like approximation with a rigid one-degree-of-freedom folding mechanism is not established, and that positive Gaussian curvature has known rigidity obstructions. This is supplied dated curation derived from the OWR record, not independently verified here; the status and cited claims require independent verification.
\item The registry marks the statement \texttt{corrected\_verified} and reports manual semantic, typographic, and LaTeX review of a meaning-preserving cleanup. This is supplied editorial metadata and does not establish mathematical correctness or resolve the formulation issues above.
\end{itemize}
\end{document}
